% !Mode:: "TeX:UTF-8"
%% 请使用 XeLaTeX 编译本文.
\documentclass{HENUMaster}
\usepackage{subfigure}
\usepackage{listings}
\usepackage{bm}
\usepackage[rm,center,compact]{titlesec}
\usepackage{caption}
\usepackage{colortbl}
\usepackage{titletoc}
\makeatletter
\renewcommand\thefigure{\@arabic\c@chapter-\@arabic\c@figure}%
\makeatother
%\makeatletter
%\renewcommand\thesubfigure{\@arabic\c@chapter--\@arabic\c@subfigure}%
%\makeatother
\makeatletter
\renewcommand\thetable{\@arabic\c@chapter-\@arabic\c@table}%
\makeatother
\titlecontents{chapter}[0pt]{\zihao{-4}\selectfont\heiti\addvspace{2pt}\filright}
              {\contentspush{\thecontentslabel\ }}
              {}{\titlerule*[8pt]{.}\contentspage}
\titlecontents{section}[8mm]{\zihao{-4}\selectfont\addvspace{2pt}\filright}
              {\contentspush{\thecontentslabel\ }}
              {}{\titlerule*[8pt]{.}\contentspage}
\titlecontents{subsection}[16mm]{\zihao{-4}\selectfont\addvspace{2pt}\filright}
              {\contentspush{\thecontentslabel\ }}
              {}{\titlerule*[8pt]{.}\contentspage}
\captionsetup[figure]{name=\heiti 图,labelsep=space}
\captionsetup[table]{name=\heiti 表,labelsep=space}
\titleformat{\chapter}{\vspace{-2cm}\centering\zihao{3}\heiti}{第\,\chinese{chapter}\,章}{1em}{}[\vspace{-0.3cm}]
\titleformat{\section}{\zihao{-3}\heiti}{\chinese{section}}{1em}{}
\titleformat{\subsection}{\zihao{4}\heiti}{(\,\chinese{subsection}\,)}{1em}{}
\titleformat{\subsubsection}{\zihao{-4}\heiti}{\arabic{subsubsection}}{1em}{}




%%%=== 参考文献=== %%%
\bibliographystyle{abbrv}        % 参考文献样式,  plain,unsrt,alpha,abbrv 等等
%%%%%%%%%%%%%%%%%%%%%%%%%%%%%%%%%%%%%%%%%%%%%%%%%%%%
\begin{document}
%%%%%%%-------------------------------------------------
\fenleihao{F000.0}  % 分类号:《中国图书资料分类法》的类号. 必填. 要根据自己的学科方向填写!!
\miji{}                % 密级
\bianhao{10475}  % 学校编号, 10475是河南大学的编号. 不用改动.

\title{河南大学硕士研究生毕业论文 \LaTeX 模板}
\Etitle{\LaTeX Template of Henan University} % 英文题目
\author{姓名}   %姓名
\StudentNumber{104750000000}   % 学号
\Eauthor{English name}    %作者英文名
\Ckind{XX学硕士}    %申请学位类别
\Csupervisor{XXX\quad 教授}        %指导教师中文名、职称
\Esupervisor{Prof. XXXX}     %指导教师英文名、职称
\Cmajor{XXXX学}                          % 专业中文名[ SMD:专业类别（领域）]
\Emajor{Major}% 专业英文名[ SMD:专业类别（领域）]
\Cspeciality{研究方向}                     % 研究方向
\Especiality{xxxx }   % 研究方向
\Schoolname{School of Economics} %学院英文名. 不确定的话, 请看一下自己学院的网页上是怎么写的. 别搞错了!
\date{二〇一七\ 年\ 六 \ 月}                    % 硕士类只写年月. 要注意和英文日期一致!!
\Edate{June, 2017}                   % 英文封面日期

%-----------------------------------------------------------------------------
\pdfbookmark[0]{封面}{title}         % 封面页加到 pdf 书签
\maketitle
%-----------------------------------------------------------------------------
% !Mode:: "TeX:UTF-8"

%%% 此部分包含: (1) 英文封面 (无需改动) ; (2) 郑重声明 (无需改动).

%%%%%%%%%%%%%%%%%%%%%%%%%%%%%
%%% -------------  英文封面 (无需改动)-------------   %%%
%%%%%%%%%%%%%%%%%%%%%%%%%%%%%
\thispagestyle{empty}
\renewcommand{\baselinestretch}{1.5}  %下文的行距
\ \vspace{1.2cm}
{\tnewroman
\begin{center}{\zihao{-2} \the\Etitle \par}\end{center}
\vspace{5cm}
\begin{center}
\zihao{-3}
A Dissertation Submitted to\\
the Graduate School of Henan University\\
in Partial Fulfillment of the Requirements\\
for the Degree of \\
Master of Economics\\
\vspace{4cm}
%\vspace*{2cm}
%\begin{center}
  %\iflib % 向图书馆提交电子文档, 使用黑白校徽.
  %\includegraphics[height=4cm]{whu.eps}       %%  黑白的. 很小, 只有 10k.
  %\else
    % \ifprint % 文档打印, 使用黑白校徽.
  %\includegraphics[height=4cm]{whu.eps}       %%  黑白的.
  %\else
  %\includegraphics[height=4cm]{whulogo.eps} %%  彩色的.
  %\fi
 % \fi
%\end{center}
\renewcommand{\baselinestretch}{2}
By\\
\emph{\the\Eauthor}\\
Supervisor:\the\Esupervisor\\

%\zihao{-2}
%\the\Schoolname\\
%{\sc Henan University}

%\vspace*{1.0cm}

Date \the\Edate

\end{center}
}
%%%%%%%--判断是否需要空白页-----------------------------
  \iflib
  \else
  \newpage
  \thispagestyle{empty}
  \cleardoublepage
  \fi
%%%%%%%-------------------------------------------------
%%%--- 加入``郑重声明'' --- %%%%%%%%%%%%%%%%%
{\pagestyle{empty}
\newpage
\vspace*{18pt}
\begin{center}{\ziju{0}\zihao{3}\heiti\em 关于学位论文独创声明和学术诚信承诺}\end{center}
\par\vspace*{20pt}
\renewcommand{\baselinestretch}{2}
{\zihao{5} \kaishu %
\begin{spacing}{1.5}
本人向河南大学提出硕士学位申请。本人郑重声明：所呈交的学位论文是本人在导师的指导下独立完成的，对所研究的课题有新的见解。据我所知，除文中特别加以说明、标注和致谢的地方外，论文中不包括其他人已经发表或撰写过的研究成果，也不包括其他人为获得任何教育、科研机构的学位或证书而使用过的材料。与我一同工作的同事对本研究所做的任何贡献均已在论文中作了明确的说明并表示了谢意。

在此本人郑重承诺：所呈交的学位论文不存在舞弊作伪行为，文责自负。
\end{spacing}
\vskip1cm

\hspace*{3cm}学位申请人（学位论文作者）签名: \raisebox{-1ex}{\underline{\makebox[4cm][c]{}}}
\vskip0.5cm
\hspace*{8.5cm}2017\ 年\hspace{0.8cm}   月 \hspace{0.8cm} 日\hspace{0.8cm}\par}

%%%--- 加入``授权书'' --- %%%%%%%%%%%%%%%%%
\vskip2cm
\begin{center}{\ziju{0}\zihao{3}\heiti\em 关于学位论文著作权使用授权书}\end{center}
\par\vspace*{18pt}
\renewcommand{\baselinestretch}{2}
{\zihao{5} \kaishu %
\begin{spacing}{1.5}
本人经河南大学审核批准授予硕士学位。作为学位论文的作者，本人完全了解并同意河南大学有关保留、使用学位论文的要求，即河南大学有权向国家图书馆、科研信息机构、数据收集机构和本校图书馆等提供学位论文（纸质文本和电子文本）以供公众检索、查阅。本人授权河南大学出于宣扬、展览学校学术发展和进行学术交流等目的，可以采取影印、缩印、扫描和拷贝等复制手段保存、汇编学位论文（纸质文本和电子文本）。

（涉及保密内容的学位论文在解密后适用本授权书）
\end{spacing}
\vskip1cm

\hspace*{3cm}学位获得者（学位论文作者）签名: \raisebox{-1ex}{\underline{\makebox[4cm][c]{}}}
\vskip0.5cm
\hspace*{8.5cm}2017\ 年\hspace{0.8cm}   月 \hspace{0.8cm} 日\hspace{0.8cm}\par
\vskip0.5cm
\hspace*{4.85cm}学位论文指导教师签名: \raisebox{-1ex}{\underline{\makebox[4cm][c]{}}}
\vskip0.5cm
\hspace*{8.5cm}2017\ 年\hspace{0.8cm}   月 \hspace{0.8cm} 日\hspace{0.8cm}\par}
%%%%%%%--判断是否需要空白页-----------------------------
  \iflib
  \else
  \newpage
  \cleardoublepage
  \fi
%%%%%%%-------------------------------------------------}
\renewcommand{\baselinestretch}{1.6}
\small\normalsize




    % 加入英文封面
\frontmatter
\pagenumbering{Roman}               % 正文之前的页码用大写罗马字母编号.

\cleardoublepage
\newpage  \pagestyle{fancy} \fancyfancy
%------------------------------------------------------------------------------
% !Mode:: "TeX:UTF-8"

%%% 说明: 此部分需要自己填写的内容:  (1) 中文摘要及关键词 (2) 英文摘要及关键词

%%%%%%%%%%%%%%%%%%%%%%%
%%% ------------ 中文摘要 ---------------%%%
%%%%%%%%%%%%%%%%%%%%%%%
\begin{cnabstract}

\par\hspace{2em}中文摘要

\end{cnabstract}
\vspace{1em}\par%\vfill

%%%--------- 关键词 -------- %%%
\cnkeywords{关键词 }

%%%%%%%%%%%%%%%%%%%%%%%


%%%%%%%%%%%%%%%%%%%%%%%
%%% ------------ 英文摘要 ---------------%%%
%%%%%%%%%%%%%%%%%%%%%%%

\begin{enabstract}
\tnewroman\zihao{5}
\par\hspace{2em} abstract



\end{enabstract}
\vspace{1em}\par%\vfill

%%%------ 英文关键词 ------- %%%
\enkeywords{\tnewroman keywords}


      % 加入中英文摘要.
%---把目录加入到书签---%%%%%%%%%%%%%%
\pdfbookmark[0]{目录}{toc}%%%%%%%%%%%%

\tableofcontents

%------------------------------------------------------------------------------
\mainmatter %% 以下是正文
\zihao{-4}
\baselineskip=24pt  % 正文行距为 20 磅
%%%%%%%%%%%%%%%%%%%%%%%%%%%%%%%%%%%%

\setlength{\parskip}{0em}

\chapter{先说重要的}

\section{具体使用步骤}
\par 本模板根据武汉大学黄正华老师的武汉大学硕士论文\LaTeX 模板修改而来，已经完全符合河南大学对硕士学位论文的格式要求。
\begin{description}

  \item[Step 1]  进入 includefile 文件夹,  打开 midmatter.tex,  backmatter.tex 这两个文档,
        分别填写 (1) 中文摘要、英文摘要, (2) 致谢.

  \item[Step 2]  打开主文档 MasterTemplate.tex, 填写题目、作者等等信息, 书写正文.

  \item[Step 3]  使用 XeLaTeX 编译. 具体见 \ref{sec-compile} 节.


\end{description}




\section{编译的方法}\label{sec-compile}

默认使用 XeLaTeX 编译, 直接生成~pdf 文件.

若另存为新文档, 请确保文档保存类型为 \verb|:UTF-8|. 当然目前很多编辑器默认文字编码为 UTF-8.
WinEdt 9.0 之后的版本都是默认保存为 UTF-8 的.



\section{文档类型选择}
文档类型有 3 种情形:

\begin{table}[ht]\centering
\begin{tabular}{ll}
\hline
   \verb|\documentclass{HENUMaster}|                 &  硕士论文 \\
   \verb|\documentclass[forprint]{HENUMaster}|    &  硕士论文打印版  \\
   \verb|\documentclass[forlib]{HENUMaster}|       &  硕士论文图书馆提交版  \\
\hline
\end{tabular}
\end{table}

另外: 专业硕士毕业论文, 请在上述情形另外加上选项 smd. 专业硕士毕业论文的封面稍有不同(中英文封面), 页眉也顺势改变了. 即
\begin{table}[ht]\centering
\begin{tabular}{ll}
\hline
   \verb|\documentclass[smd]{HENUMaster}|               &  专业硕士论文 \\
   \verb|\documentclass[smd,forprint]{HENUMaster}|    &  专业硕士论文打印版  \\
   \verb|\documentclass[smd,forlib]{HENUMaster}|       &  专业硕士论文图书馆提交版  \\
\hline
\end{tabular}
\end{table}



相关解释见接下来的两节.



\section{打印的问题}
\begin{enumerate}[i)]
  \item  硕士论文要求\colorbox{yellow}{双面打印}.
  \item  关于文档选项 forprint: 交付打印时, 建议加上选项 forprint, 以消除彩色链接文字, 避免打印字迹偏淡.
  \item  打印时留意不要缩小页面或居中. 即页面放缩方式应该是``无''(Adobe Reader XI 是选择``实际大小'').
           有可能页面放缩方式默认为``适合可打印区域'', 会导致打印为原页面大小的 $97\%$.
           文字不要居中打印, 是因为考虑到装订, 左侧的空白留得稍多一点(模板已作预留).
\end{enumerate}

\textbf{问}: {\kaishu 生成 PDF 文件时，不能去掉目录和文章的引用彩色方框，请问怎么解决？}

\textbf{答}: {\kaishu 方框表示超级链接, 只在电脑上看得见. 实际打印时, 是没有的.}





\section{提交电子文档}

向河南大学图书馆提交电子文档, 需单独编译文件: 在文档选项中添加 forlib.

原因: 图书馆要求提交的电子文档不能有空白页、彩色文字.

({\kaishu 这个要求使我在编制模板时遇到了一点问题: 这会导致电子版与纸质版的页码不一致.  论文每章的起始页默认在奇数页, 这会不可避免地出现空白页.})



\chapter{杂七杂八的话}

\section{Readme}

模板文件的结构, 如下表所示:
 \begin{table}[ht]\centering
\begin{tabular}{r|r|l}
\hline\hline
  \multicolumn{2}{l|}{MasterTemplate.tex }  &  主文档. 在其中填写正文.\\\hline
                          &frontmatter.tex&  英文封面. \hfill ({\kaishu 无需干预}) \\\cline{2-3}
 includefile 文件夹  & midmatter.tex  &  中文摘要, 英文摘要.\hfill  ({\kaishu 自行填写}) \\\cline{2-3}
                            & backmatter.tex &  致谢.\hfill  ({\kaishu 自行填写}) \\\hline
  \multicolumn{2}{l|}{figures 文件夹} &  存放图片文件.\\\hline
  \multicolumn{2}{l|}{BIBbase 文件夹} &   供 BibTeX{} 做参考文献时选用.\\
\hline
  \multicolumn{2}{l|}{HENUMaster.cls} &  定义文档格式的 class file. 不可删除.\\ \hline \hline
\end{tabular}
\end{table}
%  \footnotetext[1]{定制这个参考文献格式, 主要是希望尽可能满足河南大学硕士论文的相应格式要求, 比如要求将~article 的``年份''写在``卷号''之前.
%                   存在的问题是: 无法将~book 的``出版时间''标在``出版单位''之前. 事实上, 这个要求有悖于``国家标准~7714-87 文后参考文献著录规则''.
%                   作者名的写法, 依照了``文后参考文献著录规则'', 比如~Albert Einstein 写为~Einstein A.}
%  \footnotetext[1]{编译信息文件常常让我们的文件夹变得凌乱不堪, WinEdt窗口的``垃圾箱''按钮(Erase Output Files),
%     可以让我们方便地删除这些``垃圾''文件. 这也减少了误删文件的可能.}

无需也不要改变、移动上述文档的位置.



可能这里的文档看上去有点多. 如果不习惯用~\verb|\include{ }|~的方式加入``子文档'', 当然可以把它们合并在主文档, 成为一个文档.
({\kaishu 但是这样并不会给我们带来方便.})

利用~WinEdt~的~Project tree, 可以方便地管理这些文件:
\begin{itemize}
    \item 点击~WinEdt~窗口的~Project Tree~按键;
    \item 再点击~WinEdt~窗口的~Set Main File~按键;
\end{itemize}
接下来的管理, 已经清楚地展示在跳出的窗口中了. 再去处理其他的文件时, 还要点击~WinEdt~窗口的~Remove Main File~按键.


英文封面页加了校徽, 有黑白、彩色两种备选, 都是矢量图形, 放大不会失真. 图片文件很小,
用~MetaPost~生成的(准确的说, 是最后一道工序用到了~MetaPost), 过程有点复杂.
封面上毛体字的``河南大学''也是这样做出的.



2016 年 02 月更新: 调整为适应 TeX Live 2015 的版本.

2016 年 04 月更新: 按照研究生院新发布的要求进行改造. 主要变化: (1) 封面``河南大学''字样改为毛体题字;
(2) 郑重声明改为论文原创性声明; (3) 加了页眉; (4) 专业硕士论文封面, 供选择. 使用选项 smd 即可.

\section{字体调节}

\begin{tabular}{ll}
 \verb|\songti| & {\songti 宋体} \\
 \verb|\heiti| & {\heiti 黑体} \\
 \verb|\fangsong| & {\fangsong 仿宋} \\
 \verb|\kaishu| & {\kaishu 楷书} \\
\end{tabular}


\section{字号调节}
字号命令: \verb|\zihao|

\begin{tabular}{ll}
\verb|\zihao{0}| &\zihao{0}  初号字 English \\
\verb|\zihao{-0}|&\zihao{-0} 小初号 English \\
\verb|\zihao{1} |&\zihao{1}  一号字 English \\
\verb|\zihao{-1}|&\zihao{-1} 小一号 English \\
\verb|\zihao{2} |&\zihao{2}  二号字 English \\
\verb|\zihao{-2}|&\zihao{-2} 小二号 English \\
\verb|\zihao{3} |&\zihao{3}  三号字 English \\
\verb|\zihao{-3}|&\zihao{-3} 小三号 English  \\
\verb|\zihao{4} |&\zihao{4}  四号字 English  \\
\verb|\zihao{-4}|&\zihao{-4} 小四号 English \\
\verb|\zihao{5} |&\zihao{5}  五号字 English \\
\verb|\zihao{-5}|&\zihao{-5} 小五号 English \\
\verb|\zihao{6} |&\zihao{6}  六号字 English \\
\verb|\zihao{-6}|&\zihao{-6} 小六号 English \\
\verb|\zihao{7} |&\zihao{7}  七号字 English \\
\verb|\zihao{8} |&\zihao{8}  八号字 English \\
\end{tabular}



\section{已加入的常用宏包}

\begin{description}
%  \item[amsmath,amssymb]
  \item[cite]  参考文献引用, 得到形如 [3-7] 的样式.
  \item[color,xcolor]  支持彩色.
  \item[enumerate]  方便自由选择 enumerate 环境的编号方式. 比如

  \verb|\begin{enumerate}[(a)]| 得到形如 (a), (b), (c) 的编号.

  \verb|\begin{enumerate}[i)]| 得到形如 i), ii), iii) 的编号.

\end{description}

另外要说明的是,  itemize, enumerate, description 这三种 list 环境, 已经调节了其间距和缩进,
以符合中文书写的习惯.

\section{标点符号的问题}

建议使用半角的标点符号, 后边再键入一个空格. 特别是在英文书写中要注意此问题!

使用半角也有不足的地方, 比如顿号出不来, 需要临时切换到全角. (当然也有输入法不需要切换, 比如紫光拼音, 可以直接用拼音输入逗号.)

双引号是由两个左单引号、两个右单引号构成的: \verb|``  ''|. 左单引号在键盘上数字~1 的左边.

但是, 无论您偏向于全角或半角, 强烈建议您使用实心的句号, 只要您书写的是自然科学的文章.
原因可能是因为, 比如使用全角句号的句子结尾处的``$x$。''容易误为数学式~$x_0$(\verb|$x_0$|)吧.

\section{中英文间距问题}

自动加入间距. 不再需要在公式、英文前后加字符``\verb|~|''或空格.

\section{引用的问题}


\subsection{参考文献的引用}

参考文献的引用, 用命令~\verb|\cite{ }|. 大括号内要填入的字串, 是自命名的文献条目名.

比如, 通常我们会说:

 {\kaishu
关于此问题, 请参见文献 \cite{r2}. 作者某某还提到了某某概念\upcite{r1}.}


上文使用的源文件为:

 {\kaishu
关于此问题, 请参见文献~\verb|\cite{r2}|. 作者某某还提到了某某概念~\verb|\upcite{r1}|.
}

其中~\verb|\upcite| 是自定义命令, 使文献引用呈现为\CJKunderdot{上标形式}.

({\heiti 注意:} {\kaishu 这里文献的引用, 有时需要以上标形式出现, 有时需要作为正文文字出现, 为什么?})

另外, 要得到形如~\cite{r1,r3,r4,r5} 的参考文献连续引用, 需要用到 cite 宏包(模板已经加入),
在正文中使用~\verb|\cite{r1,r3,r4,r5}| 的引用形式即可.
或者, 连续引用的上标形式: 使用~\verb|\upcite{r1,r2,r3}|, 得到\upcite{r1,r2,r3}.

\subsection{定理和公式的引用}

\begin{theorem}[谁发现的]\label{th-abcd}
最大的正整数是~$1$.
\end{theorem}

\begin{proof}
要找到这个最大的正整数, 我们设最大的正整数为~$x$, 则~$x \geqslant 1$, 两边同时乘以~$x$, 得到
\begin{equation}\label{eq-abc}
x^2 \geqslant x.
\end{equation}
而~$x$ 是最大的正整数, 由~\eqref{eq-abc} 式得到
\[
x^2 = x.
\]
所以
\begin{equation*}
x = 1.
\end{equation*}
\end{proof}

定理~\ref{th-abcd} 是一个重大的发现.

%%%%----- 定义等环境的举例 --------
\begin{definition}[整数]
 正整数(例如 1, 2, 3)、负整数(例如 ${−1}$, $−2$, $−3$)与零(0)合起来统称为{\heiti 整数}.
\end{definition}

\begin{remark}
  整数集合在数学上通常表示为 $\mathbf{Z}$ 或 $\mathbb{Z}$, 该记号源于德语单词 Zahlen(意为``数'')的首字母.
\end{remark}

\begin{proposition}
任意两个整数相加、相减、相乘的结果, 仍然是整数.
\end{proposition}

\begin{example}
  $1+2=3$.
\end{example}

\begin{corollary}
   在整数集合内, 相加、相减、相乘运算是封闭的.
\end{corollary}

\section{图形与表格}

支持 eps, pdf, jpg 这几种常见图形格式.

再次澄清一个误会: \LaTeX{} 支持的图形格式绝非 eps 这一种. 无需特意把图片转化为 eps 格式.

用形如 \verb|\includegraphics[width=12cm]{1.jpg}| 的命令可以纳入图片.

如图 \ref{fig:1} 是一个纳入~jpg 图片的例子.

\begin{figure}[ht]
\centering
  \includegraphics[width=\textwidth]{HENU.jpg}
  \caption{一个彩色 jpg 图片的例子}
  \label{fig:1}
\end{figure}


多张图并排

\begin{figure}[!htbp]
\centering
\subfigure[\heiti\zihao{-5}1]{
\label{Fig.sub.1}
\includegraphics[width=5cm]{HENU.jpg}}
\subfigure[\heiti\zihao{-5}1]{
\label{Fig.sub.2}
\includegraphics[width=5cm]{HENU.jpg}}
\\
\subfigure[\heiti\zihao{-5}1]{
\label{Fig.sub.3}
\includegraphics[width=5cm]{HENU.jpg}}
\subfigure[\heiti\zihao{-5}11]{
\label{Fig.sub.4}
\includegraphics[width=5cm]{HENU.jpg}}
\caption{\heiti\zihao{-5}1}
\label{fig:2}
\end{figure}

表格问题, 建议使用``三线表'', 如表 \ref{tab:1}.

\begin{table}[ht]
\centering
\caption{一般三线表}
\label{tab:1}
    \begin{tabular}{c c c c c c c c c c c}
    \hline
    123 & 4  & 5  & 123 & 4 & 5123 & 4 & 5 & 123 & 4 & 5\\
    \hline
    67 & 890 & 13 & 123 & 4 & 5123 & 4 & 5 & 123 & 4 & 5\\
    67 & 890 & 13 & 123 & 4 & 5123 & 4 & 5 & 123 & 4 & 5\\
    67 & 890 & 13 & 123 & 4 & 5123 & 4 & 5 & 123 & 4 & 5\\
    \hline
    \end{tabular}
\end{table}

%\begin{table}[h]
%  \centering
%  \caption{调整线宽的三线表}\label{tab:2}
%    \begin{tabular}{c c c c c c c c c c c}
%    \hlinewd{0pt}
%    姓名: &&&&年龄: &&&& 性别:\\
%    \hlinewd{1.5pt}
%    123 & 4  & 5  & 123 & 4 & 5123 & 4 & 5 & 123 & 4 & 5\\
%    \hlinewd{1pt}
%    67 & 890 & 13 & 123 & 4 & 5123 & 4 & 5 & 123 & 4 & 5\\
%    67 & 890 & 13 & 123 & 4 & 5123 & 4 & 5 & 123 & 4 & 5\\
%    67 & 890 & 13 & 123 & 4 & 5123 & 4 & 5 & 123 & 4 & 5\\
%    \hlinewd{1.5pt}
%    \end{tabular}
%\end{table}

%\section{论文电子版提交的格式问题}
%
%《河南大学博、硕士研究生学位论文缴送办法》要求``电子版应采用~Word97 或~Word2000 格式(DOC)编辑''\footnote{那已经是~2005 年的事情了.}.
%关于这个重要的问题, 我已致信河南大学图书馆总馆信息服务中心袁晓川老师,
%得到的答复是: ``用~\LaTeX{} 软件排版的论文可直接提交~PDF 格式.''







\chapter{其他事项}
以下是广告时间, 插播一段广告:
\begin{itemize}
    \item 插图的制作, 建议 pgf.
          pgf 的长处是源文件直接植入~\TeX~文档, 管理起来非常方便.
    这里有我写的一个关于初次使用~pgf~的帖子:\\    \url{http://bbs.ctex.org/forum.php?mod=viewthread&tid=30480}.
    \item 生成参考文献, 建议使用~BibTeX.  这里有我写的一个文档: \\
    \url{http://bbs.ctex.org/forum.php?mod=viewthread&tid=26056}.

      {\kaishu 使用 BibTeX{} 做参考文献时,
      借助 EndNote 或者 NoteExpress, 可以非常漂亮简单地解决 bib 文件的录入问题.
      NoteExpress 在校图书馆网站有正版软件提供下载.
      当然 EndNote 本身就是 Thomson Corporation 推出的(和 SCI 搜索引擎是同一家公司),
      和多个重要文献搜索引擎有良好的功能配合.

      Google 学术搜索也提供了文献的 bib 格式.
      录入参考文献时, 偶尔用一用 Google 学术搜索, 还可以核查或减少录入的错误, 并减少录入的工作量.}
     \item 幻灯片的制作, 建议使用~Beamer. 这里有我写的一个模板, 谨供参考:\\
    \url{http://bbs.ctex.org/forum.php?mod=viewthread&tid=27695}.
\end{itemize}

%%%=====================================================================%%%

%%%=== 参考文献 ========%%%
\cleardoublepage\phantomsection
\addcontentsline{toc}{chapter}{参考文献}
\begin{thebibliography}{000}\zihao{5}

  \bibitem{r1} 作者, 文章题目, 期刊名, 年份(期数): 起止页码

  \bibitem{r2} 作者, 书名, 年份, 版次, 出版地: 出版单位, 起止页码

  \bibitem{r3} 邓建松等, 《\LaTeXe~科技排版指南》, 科学出版社

  \bibitem{r4} 吴凌云, 《CTeX~FAQ (常见问题集)》, \textit{Version~0.4}, June 21, 2004

  \bibitem{r5} Herbert Vo\ss, Mathmode, \url{http://www.tex.ac.uk/ctan/info/math/voss/mathmode/Mathmode.pdf}.

\end{thebibliography}



\backmatter
% !Mode:: "TeX:UTF-8"
%%%%%%%%%%%%%%%%%%%%%%%
%%% --------------- 致谢 ------------- - %%%
%%%%%%%%%%%%%%%%%%%%%%%
\acknowledgement
\zihao{-4}

\par \qquad 致谢！


\vskip3em
\rightline{
\newcommand{\xiaoqi}{\fontsize{12.5pt}{12pt}\selectfont}
\xiaoqi
\begin{tabular}{cc}
\centering
XXX&\\
2017年5月&\\
于河南大学金明园&
\end{tabular}\qquad}



 %%%致谢, 河南大学学位论文使用授权协议书．
\cleardoublepage
\end{document}



